\documentclass{ximera}
\author{Jeffrey Kuan}
\title{Math 1050 Notes: Section 2.1}
\license{CC: 0}

\begin{document}

\begin{abstract}
    Section 2.1 --- Evaluating Variable Expressions.
    Vocabulary, worked examples, then practice problems with answers.
\end{abstract}

\maketitle

\section*{Vocabulary}

\begin{itemize}
  \item \textbf{Expression} --- not a complete sentence, for example $5x+3$.
  \item \textbf{Equation} --- a complete sentence, containing an equals sign, for example
        $5x+3=15$.
  \item \textbf{Variable} --- something that varies; its value can change. Usually a
        letter.
  \item \textbf{Coefficient} --- tells how many of a variable, as the 5 in $5x$ or the 3
        in $3a$.
  \item \textbf{Constant} --- something that does not change; a number with no variable.
  \item \textbf{Term} --- the parts separated by $+$ or $-$. For example $x+3$ has two
        terms.
\end{itemize}

\begin{problem}
\textbf{1.} How many terms are in each variable expression?

(a) $2x^{2}+3x-9$ \qquad \textbf{Answer:} $\answer{3}$

(b) $7x^{2}y+4xy-6xy^{2}+1$ \qquad \textbf{Answer:} $\answer{4}$

(c) $5xyz$ \qquad \textbf{Answer:} $\answer{1}$
\end{problem}

\section*{Evaluate when $a=2$, $b=3$, $c=-4$}

\begin{problem}
\textbf{2.} $a-2c$ \qquad \textbf{Answer:} $\answer{10}$

\begin{hint}
\textbf{Solution.} $2-2(-4)=2+8=10$.
\end{hint}
\end{problem}

\begin{problem}
\textbf{3.} $-3c+4$ \qquad \textbf{Answer:} $\answer{16}$

\begin{hint}
\textbf{Solution.} $-3(-4)+4=12+4=16$.
\end{hint}
\end{problem}

\begin{problem}
\textbf{4.} $2a-c$ \qquad \textbf{Answer:} $\answer{8}$

\begin{hint}
\textbf{Solution.} $2(2)-(-4)=4+4=8$.
\end{hint}
\end{problem}

\begin{problem}
\textbf{5.} $ab+abc$ \qquad \textbf{Answer:} $\answer{-18}$

\begin{hint}
\textbf{Solution.} $(2)(3)+(2)(3)(-4)=6-24=-18$.
\end{hint}
\end{problem}

\begin{problem}
\textbf{6.} $\dfrac{a-c}{c}$ \qquad \textbf{Answer:} $\answer{-\frac{3}{2}}$

\begin{hint}
\textbf{Solution.} $\dfrac{2-(-4)}{-4}=\dfrac{6}{-4}=-\dfrac{3}{2}$.
\end{hint}
\end{problem}

\begin{problem}
\textbf{7.} $-\dfrac{1}{2}c-\dfrac{2}{3}b$ \qquad \textbf{Answer:} $\answer{0}$

\begin{hint}
\textbf{Solution.} $-\dfrac{1}{2}(-4)-\dfrac{2}{3}(3)=2-2=0$.
\end{hint}
\end{problem}

\begin{problem}
\textbf{8.} $(c-a)^{2} \div b$ \qquad \textbf{Answer:} $\answer{12}$

\begin{hint}
\textbf{Solution.} $(-4-2)^{2}\div 3=(-6)^{2}\div 3=36\div 3=12$.
\end{hint}
\end{problem}

\begin{problem}
\textbf{9.} $\dfrac{a}{-c} \div \dfrac{b}{a}$ \qquad \textbf{Answer:} $\answer{\frac{1}{3}}$

\begin{hint}
\textbf{Solution.} $\dfrac{2}{4}\div\dfrac{3}{2}=\dfrac{1}{2}\cdot\dfrac{2}{3}=\dfrac{1}{3}$.
\end{hint}
\end{problem}

\section*{Volume. Round to the nearest tenth, using $3.14$ for $\pi$}

\begin{problem}
\textbf{10.} Find the volume of a sphere with radius $8.5$ cm.
\qquad \textbf{Answer:} $\answer{2571.1}$

\begin{hint}
\textbf{Solution.} $V=\dfrac{4}{3}\pi r^{3}=\dfrac{4}{3}(3.14)(8.5)^{3}=2571.1$ cm$^{3}$.
\end{hint}
\end{problem}

\begin{problem}
\textbf{11.} Find the volume of a right circular cylinder with radius $1.5$ in and height
$5.5$ in. \qquad \textbf{Answer:} $\answer{38.9}$

\begin{hint}
\textbf{Solution.} $V=\pi r^{2}h=(3.14)(1.5)^{2}(5.5)=38.9$ in$^{3}$.
\end{hint}
\end{problem}

\section*{Your turn}

\begin{problem}
\textbf{1.} Evaluate each expression when $a=-2$, $b=4$, $c=-1$, $d=-3$.

(a) $\dfrac{b+c}{d}$ \qquad \textbf{Answer:} $\answer{-1}$

(b) $\dfrac{b-d}{c-d}$ \qquad \textbf{Answer:} $\answer{\frac{7}{2}}$

(c) $b^{2}-2b+4$ \qquad \textbf{Answer:} $\answer{12}$

(d) $\dfrac{2}{3}d^{2}-\dfrac{3}{8}b^{2}$ \qquad \textbf{Answer:} $\answer{0}$

\begin{hint}
\textbf{Solutions.} (a) $\dfrac{4-1}{-3}=-1$. \quad
(b) $\dfrac{4+3}{-1+3}=\dfrac{7}{2}$. \quad
(c) $16-8+4=12$. \quad
(d) $\dfrac{2}{3}(9)-\dfrac{3}{8}(16)=6-6=0$.
\end{hint}
\end{problem}

\begin{problem}
\textbf{2.} Find the volume of a sphere with radius $3$ in. Round to the nearest tenth.
\qquad \textbf{Answer:} $\answer{113.0}$

\begin{hint}
\textbf{Solution.} $V=\dfrac{4}{3}(3.14)(3)^{3}=113.0$ in$^{3}$.
\end{hint}
\end{problem}

\end{document}
